\documentclass[11pt,a4paper]{article}
\usepackage[utf8]{inputenc}
\usepackage[T1]{fontenc}
\usepackage{lmodern}
\usepackage[margin=2.4cm]{geometry}
\usepackage{amsmath,amssymb}
\usepackage{booktabs}
\usepackage{graphicx}
\usepackage{xcolor}
\usepackage[colorlinks=true,linkcolor=blue,citecolor=blue,urlcolor=blue]{hyperref}

\newcommand{\Hc}[1]{H_{c#1}}
\newcommand{\tGL}{\tau_{\mathrm{GL}}}
\newcommand{\code}[1]{\texttt{#1}}

\title{juliaTDGL: a guide to the three-dimensional two-band\\
time-dependent Ginzburg--Landau solver}
\author{J.~A.~Robles-Calder\'on\\
\small Departamento de F\'isica, Universidad Nacional de Colombia, Bogot\'a}
\date{\today}

\begin{document}
\sloppy
\maketitle

\begin{abstract}
\noindent
This guide describes what the solver computes and how: the two-band
model, its discretisation on a staggered lattice with link variables, the
time step, the injection of a transport current, the voltage observable,
the convergence of a current ladder, the optional scalar potential and the
optional terms (mixed-gradient coupling, heating, Joule heating, bath
temperature, charge imbalance), the geometry, and the protocol for
reporting a diode efficiency. It ends with how to run the code, what it
writes, and what each test checks.
\end{abstract}

\tableofcontents

% ======================================================================
\section{What the solver computes}

The solver integrates the time-dependent Ginzburg--Landau (TDGL) equations
of a superconductor with two order parameters, $\Psi_\sigma$ and $\Psi_\pi$,
coupled by a Josephson term, together with the vector potential
$\mathbf A$, on a three-dimensional rectangular box surrounded by a buffer,
under an applied field and a transport current. It was written to study
the geometric superconducting diode: a sample whose shape breaks the
mirror symmetry along the direction of vortex motion carries different
critical currents in the two directions.

The output of a run is a current--voltage characteristic for both
polarities, the vortex content and magnetisation at every current, and the
full state at chosen points, from which the critical currents
$I_c^\pm$ and the efficiency
\begin{equation}
  \eta = \frac{I_c^{+}-|I_c^{-}|}{I_c^{+}+|I_c^{-}|}
\end{equation}
follow.

% ======================================================================
\section{The model}
\label{sec:model}

\subsection{Free energy and units}

The free energy is the two-component Ginzburg--Landau functional with an
interband Josephson coupling,
\begin{equation}
  \mathcal{F} = \sum_{a=\sigma,\pi}\Bigl[
      \alpha_a|\Psi_a|^2 + \frac{\beta_a}{2}|\Psi_a|^4
    + \frac{1}{2m_a}\bigl|(-i\nabla-\mathbf{A})\Psi_a\bigr|^2 \Bigr]
  - \gamma_{12}\bigl(\Psi_\sigma^{*}\Psi_\pi+\mathrm{c.c.}\bigr)
  + \tfrac{1}{2}|\nabla\times\mathbf{A}|^2 .
\end{equation}
Lengths are measured in the $\sigma$-band coherence length $\xi_\sigma$,
fields in $\Hc2\equiv\Hc2^\sigma$, time in $\tGL$, order parameters in the
$\sigma$-band equilibrium amplitude at $T=0$ and currents in the natural
unit $J_0$ of the equations. With this scaling the remaining parameters are
\begin{center}
\begin{tabular}{lll}
\toprule
symbol & \code{Config} field & meaning \\
\midrule
$\kappa$ & \code{kappa\_s} & Ginzburg--Landau parameter, $\lambda/\xi_\sigma$ \\
$a_{2n}$ & \code{a2n} & equilibrium density of the $\pi$ band at $T=0$ (default $0.59$) \\
$m_r$ & \code{m\_ratio} & mass ratio $m_\pi/m_\sigma$ (set from \code{xi\_ratio} when given) \\
$R$ & \code{xi\_ratio} & coherence-length ratio $\xi_\pi/\xi_\sigma$ \\
$\gamma_{12}$ & \code{gamma12} & interband Josephson coupling \\
$\Gamma_r$ & \code{gamma\_ratio} & relaxation time of $\pi$ relative to $\sigma$ \\
$T$ & \code{T\_red} & reduced temperature \\
$H$ & \code{H\_const} & applied field along $z$, in $\Hc2$ \\
\bottomrule
\end{tabular}
\end{center}

\subsection{Equations of motion}

In the gauge $\Phi=0$ (the default) the dynamics is relaxational:
\begin{align}
  \partial_t\Psi_\sigma &= (\nabla-i\mathbf A)^2\Psi_\sigma
     + \bigl(f\alpha_\sigma-|\Psi_\sigma|^2\bigr)\Psi_\sigma + \gamma_{12}\Psi_\pi ,\\
  \Gamma_r\,\partial_t\Psi_\pi &= \frac{1}{m_r}(\nabla-i\mathbf A)^2\Psi_\pi
     + \bigl(f\alpha_\pi-|\Psi_\pi|^2\bigr)\Psi_\pi + \gamma_{12}\Psi_\sigma ,\\
  \partial_t\mathbf A &= \mathbf J_s - \kappa^2\nabla\times\nabla\times\mathbf A ,
  \qquad
  \mathbf J_s=\mathrm{Im}[\Psi_\sigma^*(\nabla-i\mathbf A)\Psi_\sigma]
     +\frac{1}{m_r}\mathrm{Im}[\Psi_\pi^*(\nabla-i\mathbf A)\Psi_\pi] ,
\end{align}
with $\alpha_\sigma=1-T$ and $\alpha_\pi=a_{2n}-T$. The normal conductivity
is one, which fixes the voltage scale. $\Gamma_r$ divides the whole
right-hand side of the $\pi$ equation: it sets a relaxation time and does not
move any static quantity.

\paragraph{Temperature.} By default $T$ enters only the $\pi$ condensation
term and the field equation carries $(1-T)\mathbf J_s$, the convention of
the methods paper. Both conventions
coincide at $T=0$. Setting \code{T\_physical=true} gives the equations
above, with $T$ felt by both bands and no factor on $\mathbf J_s$; the
coherence-length ratio is then fixed at its $T=0$ value. At $T=0.3$ the
uniform equilibrium is $|\Psi_\sigma|^2=1-T$ and $|\Psi_\pi|^2\simeq a_{2n}-T$.

\subsection{Coherence lengths, critical fields and depairing}

Linearising about $\Psi=0$ with the bands decoupled,
\begin{equation}
  R=\frac{\xi_\pi}{\xi_\sigma}=\frac{1}{\sqrt{m_r a_{2n}}},\qquad
  \Hc2^\pi=m_ra_{2n}\Hc2=\frac{\Hc2}{R^2}.
\end{equation}
The solver takes $R$ as input and sets $m_r=1/(a_{2n}R^2)$; it prints the
ratio it realises. Raising $R$ lowers the $\pi$-band upper critical field as
$R^{-2}$, so a $\pi$ band with a long coherence length is normal at the
fields at which vortices enter: check $\langle|\Psi_\pi|^2\rangle$ before
reading any band-resolved number. At finite $\gamma_{12}$ there is one
$\Hc2$, the largest root of $(1-H)(a_{2n}-H/m_r)=\gamma_{12}^2$.

The superfluid stiffness of the pair is $\rho_s=1+a_{2n}/m_r=1+a_{2n}^2R^2$
and the effective penetration depth $\lambda_{\rm ef}=\kappa/\sqrt{\rho_s}$.
$J_0$ is not the depairing current: a uniform superflow gives
$j(q)=q[(1+a_{2n}/m_r)-q^2(1+1/m_r^2)]$, with maximum $0.546\,J_0$ at
$R=1.30$ and $2/3\sqrt3=0.385\,J_0$ for one band.

\subsection{Geometry as a condensation mask}

The sample is always the full box. The geometry enters through a mask
$f(\mathbf r)$ multiplying the condensation coefficients: $f=1$ in the
bulk, $0<f<1$ for weakened and $f>1$ for strengthened condensation. Where
$f=0$ the condensation energy vanishes and the material is locally at its
critical temperature: a notch is neither a vacuum cut nor a normal metal,
which would need $f<0$. The local upper critical fields
are $f\Hc2$ for $\sigma$ and $f\Hc2/R^2$ for $\pi$: any region with $0<f<H$
is normal and can shunt the sample (the solver warns at start).

% ======================================================================
\section{Discretisation}

\subsection{The lattice}

The mesh is staggered. $\Psi_\sigma$ and $\Psi_\pi$ live on nodes; the
vector potential lives on links as the phase
$\theta_x=\int A_x\,\mathrm dx$ of the link variable
$U_x=e^{-i\theta_x}$; the field $\mathbf h=\nabla\times\mathbf A$ lives on
plaquettes. The covariant derivative is $\Psi_{i+1}U_x(i)-\Psi_i$, which is
exactly gauge invariant on the lattice. The sample occupies
$N_x\times N_y\times N_z$ cells of size $\mathrm dx$ inside a box padded
by \code{buf} buffer cells on every face (\code{bufx}, \code{bufy},
\code{bufz} set them per axis). In the buffer the supercurrent vanishes and
$\partial_t\mathbf A=-\kappa^2\nabla\times\nabla\times\mathbf A$: it is a
normal metal of unit conductivity, not vacuum, and carries normal current
wherever the electric field is nonzero.

\subsection{Boundary conditions}

The applied field is imposed on $h_z$ at the outer faces of the
buffer; inside, the field is whatever $\nabla\times\mathbf A$ gives, so
screening, magnetisation and vortex entry are computed rather than
prescribed. The order parameters satisfy
$\hat n\cdot(\nabla-i\mathbf A)\Psi=0$ on the six faces of the sample.
Each link type has its own mask: a link is live when both of its nodes
belong to the sample. A single mask for all three directions would include
the low-index transverse rows and exclude the high-index ones, biasing the
field equation and the voltmeter towards one edge.

The buffer is a physical parameter, not only a numerical one: the
field is clamped at its outer face, inside the region where screening
currents live when the padding is thinner than $\lambda$. On the notched
production sample each doubling of the padding admits one more vortex at
zero current. Quote absolute numbers together with the padding.

\subsection{Time stepping}

Integration is explicit. One substep applies, in this order:
\begin{center}
\begin{tabular}{cl}
\toprule
block & operation \\
\midrule
{[A]} & update $\Psi_\sigma,\Psi_\pi$ (kinetic, condensation and coupling terms) \\
{[B]} & supercurrent on every link, from $\Psi$ and the links of the previous step \\
{[C]} & tentative increment of $\theta$: supercurrent minus $\kappa^2\nabla\times\mathbf h$ \\
{[C2]} & (with \code{use\_phi}) projection by the scalar potential \\
{[E]} & links $U=e^{-i\theta}$ \\
{[F]} & field $\mathbf h=\nabla\times\theta$ on plaquettes \\
{[G]} & applied field and self-field of the transport current on the boundary \\
{[H]} & Neumann condition for $\Psi$ on the sample faces \\
{[I]} & advance: the new state becomes the old one (array swaps, no copies) \\
\bottomrule
\end{tabular}
\end{center}
The time step is the minimum over the diffusion branches of both bands and
the field equation,
\begin{equation}
  \mathrm dt=0.99\min\Bigl[\frac{\eta\,\mathrm dl}{4\sqrt{1+\gamma^2}},
  \frac{\eta\,\mathrm dl\,m_r\Gamma_r}{4\sqrt{1+\gamma^2}},
  \frac{\mathrm dl}{4\kappa_\sigma^2},\frac{\mathrm dl}{4\kappa_\pi^2},
  \frac{\nu\,\mathrm dl}{4\zeta^2}\Bigr],\qquad
  \mathrm dl=\frac{2}{\mathrm dx^{-2}+\mathrm dy^{-2}+\mathrm dz^{-2}},
\end{equation}
with $\eta=5.79$, $\nu=0.03$, $\zeta=0.5$, $\gamma=10$. At $\kappa=2$ and
$\mathrm dx=0.5$ the flat branch binds and $\mathrm dt=4.95\times10^{-3}\,\tGL$;
above $\kappa\simeq2.9$ the field branch $\mathrm dl/4\kappa^2$ takes over,
so a run at $\kappa=25$ takes about $160$ times more steps per unit of
physical time than one at $\kappa=2$. Optional diffusive terms add the
condition $\mathrm dt\le0.15\,\mathrm dx^2/D$. A \emph{cycle} is
\code{lMax}$=40$ substeps; convergence is tested once per cycle. This is the
scheme of release 1.0.0 and the library default, \code{Config(:legacy)}.

\paragraph{The production scheme.} \code{Config(:prod)} (release 1.1) removes
the three bounds that do not belong to the equations integrated here:
\begin{itemize}
\item the term $\kappa^2\nabla\times\nabla\times\theta$ of block [C] is taken
  at the new time level (\code{maxwell\_implicit}); the matrix is constant and
  separable for $\mathrm dx=\mathrm dy=\mathrm dz$, so one solve per step
  costs a few transforms per axis, and the bound $\mathrm dl/4\kappa^2$
  disappears;
\item the diffusion bounds are those of explicit Euler on diffusion plus
  reaction, $2/(8/\mathrm dl+2)$ for $\sigma$ (\code{cfl\_legacy=false}),
  without the factor $1/\sqrt{1+\gamma^2}$ and the fixed $0.03\,\mathrm dl$;
\item the $\pi$ band, whose bound is the smallest at $\xi_\pi/\xi_\sigma=4$,
  advances in substeps with $\sigma$ frozen in the coupling
  (\code{pi\_substeps=0}, automatic);
\item the step is half the bound of $\sigma$ (\code{dt\_scale=0.5}): at
  $0.99$ of it the stiffest mode changes sign every step and a flowing state
  at $\kappa=25$ loses $7\,\%$ of its voltage.
\end{itemize}
It also closes a resistive point once it carries ten crossings per line
after \code{t\_min\_GL\_res} (\code{conv\_by\_time}), reads \code{eps\_conv}
per cycle of the 1.0.0 step so that the fixed-point criterion is the same in
physical time (\code{eps\_conv\_legacy}), and uses one BLAS thread. Against
the explicit scheme it gives the same onset brackets at $\kappa=2$ and
$\kappa=25$ and flowing voltages within $0.16\,\%$, at about three times the
speed per unit time at $\kappa=2$ and about a hundred times at $\kappa=25$.
The step cannot grow much further: in strong flow at $\kappa=25$ the voltage
is $0.16\,\%$ off at $\mathrm dt=0.03$ and $3\,\%$ off at $0.04$, because the
field and the links are updated once per step during a phase slip. With heat,
a bath temperature or drag the $\pi$ band takes no substeps and only the
curl-curl term is implicit, about thirty times faster than the explicit scheme
at $\kappa=25$; heated and dragged flowing and collapsed states at $\kappa=25$
agree with the explicit ones within $0.25\,\%$ in $V$. The scalar potential
runs in the explicit scheme only.

\subsection{Symmetry}

With $\mathbf H\parallel\hat z$ and $\mathbf J\parallel\hat x$ the Lorentz
force is along $y$. Complex conjugation, $\Psi\to\Psi^*$ and
$\mathbf A\to-\mathbf A$, composed with the mirror $M_y$ maps $(+J,+H)$ onto
$(-J,+H)$. A geometry invariant under $M_y$ therefore gives
$V(-J)=-V(+J)$ and $\eta=0$ to machine precision; a diode needs a shape that
breaks $M_y$, and an asymmetry along $x$ alone does not. On the lattice the
mirror of the interior nodes is $j\mapsto(y_1+y_2+1)-j$. The functions
\code{check\_my\_symmetry}, \code{check\_mx\_symmetry} and
\code{check\_rz\_symmetry} count the cells that break each symmetry and are
printed at the start of every run.

% ======================================================================
\section{Driving and measuring}

\subsection{Current injection}

With $\Phi=0$, a current imposed as a uniform phase on the links of a
contact face is a gauge transformation: it cancels in the lattice curl and
in the Neumann condition, and the solver integrates the zero-current
problem whatever value is imposed. The solver therefore injects the
current through its own field: for a slab of thickness $d$ carrying
$\mathbf J\parallel\hat x$, $h_y=\mp\mathcal G Jd/2\kappa^2$ on the two $z$
faces (\code{inject\_mode="selffield"}, gain \code{inject\_gain}$=\mathcal G$).
This field is the curl itself, cannot be gauged away, and is odd in $J$.
It adds in quadrature to the applied field; keep the total below $\Hc2$.
The legacy phase clamp (\code{inject\_mode="phase"}) is kept only to
reproduce old runs.

\subsection{The voltage}

The voltage is read from the Josephson relation on the gauge-invariant
phase between two measuring planes inset \code{delta\_v} cells from the
contact faces,
\begin{equation}
  V_\sigma=-\frac{\mathrm d}{\mathrm dt}\Bigl\langle\sum_{\rm path}
  \arg(\Psi_{\sigma,i}^*U_{ij}\Psi_{\sigma,j})\Bigr\rangle_{\rm lines}.
\end{equation}
Each transverse line is unwrapped in time separately and only then
averaged, so every $2\pi$ slip is counted: $V=2\pi N_{\rm cross}/(n_{\rm lin}\Delta t)$.
The same construction on $\Psi_\pi$ gives $V_\pi$. The voltage is the slope
of a least-squares fit over an averaging window. The plane-to-plane
difference of $A_x$ is a static quantity and is not a voltage; it is
printed as a diagnostic (\code{Vest}). Unwrapping is faithful while every
line advances by less than $\pi$ per cycle. The solver records the largest
per-line increment of each point (\code{dmax\_arr}), which in developed flux
flow is an order of magnitude above the mean rate $|V|\,\mathrm dt_{\rm cycle}$;
when it approaches $\pi$ the remedy is to sample the phase more often
(\code{unwrap\_every}), not a smaller time step. A wrapped increment cannot
exceed $\pi$; the number of increments above 2\,rad (\code{nbig\_arr})
bounds the slips that may have been misattributed, which happens when a
core passes close to a link. It also reports the voltage on the lines whose
whole path has $f=1$ (\code{Vf1\_arr}), which tells whether slips counted
inside a notch, where $\arg(\Psi^*U\Psi)$ is poorly conditioned, contribute
to $V$, and the rate of the line integral of $\mathbf A$ over the window
(\code{VA\_arr}), which has no core to misattribute but is not gauge
invariant: with $\Phi=0$ the condensate phase between the planes also
drifts, and the two voltages differ. For a point that closes as a fixed
point, \code{V\_arr} is the slope of the last \code{V\_win} cycles, a
short window ($0.1\,\tGL$ at $\kappa=25$); \code{Vh\_arr} and
\code{Vh2\_arr} give the slope of the $\sigma$ and $\pi$ phases over the
second half of every point, which is what the verdict rule needs.
\code{keep\_rung\_states=true} keeps
the state of every closed point in \code{restart/rungs/}, for analysis after
the run.

\subsection{Vortex counting and magnetisation}

The vortex count is the fluxoid $\oint D=2\pi n-\Phi_{\rm enc}$ summed over
the plaquettes of the mid-thickness plane, excluding plaquettes where
$f\simeq0$. The mask is the geometry, never the amplitude: masking by
$|\Psi|$ removes the cores being counted. The magnetisation is
$M=\langle h_z\rangle_{\rm sample}-H$.

\subsection{A current ladder}

A run sweeps a list of currents (\code{J\_list}), first $+J$ and then $-J$,
each polarity starting from the same zero-current seed. Each rung starts
from the converged state of the previous one. A rung closes when either
\begin{itemize}
\item the state is a fixed point, $\max|\Delta|\Psi||<$\code{eps\_conv}
  per cycle (conv \code{F}); or
\item the state is a limit cycle --- flux flow --- and the windowed mean of
  $V$ is stationary (conv \code{V}),
\end{itemize}
and in both cases only after a floor of physical time, \code{t\_min\_GL}
for a static point and \code{t\_min\_GL\_res} once it is resistive. The
floors are in units of $\tGL$, not of cycles, because the length of a cycle
moves with $\mathrm dt$. The tolerance \code{eps\_conv} is per cycle;
\code{eps\_conv\_rate} sets it per unit time instead, so that it does not
change with $\mathrm dt$. Near onset a flowing state can spend most of each
period nearly frozen, and no local tolerance tells it from a fixed point:
the verdict rule below does. A rung that reaches \code{mMax} cycles is
marked \code{*} and is not converged. \code{Config(:paper)} sets the
protocol of the methods paper (tolerance, floors and ceiling).

After each rung the state of that polarity is written to
\code{restart/state\_Jpos.jld2} or \code{state\_Jneg.jld2} (compressed JLD2, exact;
\code{checkpoint\_format="octave"} writes the text \code{.mat} of the earlier releases); a run restarted with \code{restart\_pos/restart\_neg} continues at the next
rung, and \code{read\_state} and \code{load\_checkpoint} open either format, so older
runs resume too. Binary MATLAB files need \code{using MAT} (an optional extension). A file
\code{STOP.txt} in the run folder stops the run after the current rung.
\code{guard\_hours} sets a wall-clock budget, and \code{guard\_psi} stops a
run whose $|\Psi|$ diverges (usually a time step too large for $\kappa$).

% ======================================================================
\section{The scalar potential}

With \code{use\_phi=true} the solver leaves the gauge $\Phi=0$. Each
substep computes the tentative link increment $\delta$ of block [C],
solves
\begin{equation}
  D(G\Phi) = D\delta ,
\end{equation}
with $G$ the lattice gradient from nodes to links and $D$ its adjoint, and
corrects $\delta'=\delta-G\Phi$, so that $D\delta'=0$ to machine precision
(Coulomb gauge). The order parameters gain $-i\Phi\Psi$ with the time step
$\mathrm dt$ in both bands. $D\circ G$ is the seven-point Laplacian; on the
uniform box it is separable and is diagonalised once by three
one-dimensional eigenproblems (\code{phi\_method=:separable}, the default),
which is an exact direct solve; a sparse LU (\code{:lu}) gives the same
result. With the separable solver a cycle costs $1.55$ times a cycle in
the gauge $\Phi=0$.

$\Phi$ is a gauge choice: $|\Psi|$, $\mathbf h$ and the Josephson voltage
are the same with $\Phi=0$, with Dirichlet and with Neumann conditions on
$\Phi$ (\code{phi\_bc}); only the split of the voltage between the drop of
$\Phi$, printed as $V_\Phi$, and the drift of the phase changes. Neumann is
the natural condition here: there are no electrodes.

Two extensions change the physics:
\begin{itemize}
\item \textbf{Thomas--Fermi screening} (\code{lambda\_TF}$>0$): the operator
  becomes $D\circ G-\lambda_{\rm TF}^{-2}$, so the potential is screened
  over $\lambda_{\rm TF}$ and the gauge equivalence is broken.
\item \textbf{Charge imbalance} (\code{tau\_Q}$>0$, with Thomas--Fermi): a
  field $Q$ relaxes towards the projection residual,
  $\tau_Q\,\mathrm dQ/\mathrm dt+Q=D\delta'$ (implicit Euler), and the
  order parameters rotate with $\mu=\Phi+Q$. Without Thomas--Fermi the
  residual is zero and $Q$ does nothing. With \code{use\_Qss} the residual
  also diffuses, $\tau_Q(\partial_tQ-D_Q\nabla^2Q)+Q=D\delta'$.
\end{itemize}

% ======================================================================
\section{Optional terms}

All are off by default. Use one at a time against a control run with the
same protocol; with two on, a change cannot be attributed. The coefficients
are model values that switch a term on, not material constants.

\begin{description}
\item[Mixed-gradient coupling] (\code{use\_drag}, \code{nu\_drag}$=\nu$).
  Adds $\nu\,\mathrm{Re}[(\mathbf D\Psi_\sigma)^*\cdot\mathbf D\Psi_\pi]$ to the
  functional, with the corresponding terms in both order-parameter
  equations and in the supercurrent. The $\pi$-side term carries
  $\nu/\Gamma_r$ and no extra $1/m_r$, which keeps the coupling variational.
\item[Heating] (\code{use\_heat}, \code{D\_T}, \code{eta\_J},
  \code{gamma\_bath}). A temperature field obeys
  $\partial_t\theta=D_T\nabla^2\theta+\eta_J s-\gamma(\theta-\theta_{\rm bath})$,
  with $\alpha_\sigma=(1-\theta)f$ and $\alpha_\pi=(a_{2n}-\theta)f$ and the
  source $s=|\partial_t\Psi_\sigma|^2+\Gamma_r|\partial_t\Psi_\pi|^2$. A rung
  whose temperature exceeds $0.45$ is cut (conv \code{T}) and ends that
  polarity.
\item[Joule heating] (\code{use\_joule}). Adds the normal-current term
  $|\mathbf E|^2$ to $s$. Requires heating and $\Phi=0$.
\item[Contact margin] (\code{heat\_margin}). Excludes from $s$ the cells
  next to each injection face. There the injected normal current converts
  into supercurrent within about one coherence length, with an electric
  field and a rotating phase even in the Meissner state; that heat belongs
  to the contact, not to the sample.
\item[Bath temperature] (\code{T\_physical}, \code{T\_red}). See
  Sec.~\ref{sec:model}.
\end{description}

% ======================================================================
\section{Geometry}

\code{GEOM\_MODE} selects a shape from the built-in catalogue (\code{SIM},
the clean box; \code{WEDGE}, a sawtooth notch on the $y_1$ face;
\code{STAGGER}, \code{FUNNEL}, \code{TRI\_Y1}, \ldots); the suffix
\code{\_REF} gives the exact $M_y$ mirror, which is the sign control. Any
shape can also be read from a file (\code{GEOM\_MODE="CUSTOM"} or
\code{"CUSTOM\_REF"}, \code{geom\_file}): a TOML list of regions in units of
$\xi_\sigma$, painted in order,
\begin{verbatim}
[[shape]]
type = "triangle"     # circle rect triangle ngon polygon sphere pixels voxels
x = 10.0              # centre; origin at the corner (x1, y1, z1) of the sample
y = 0.0
w = 4.0
h = 3.0
skew = 1.0            # apex offset along the base
rot = 0.0             # degrees about z
z = "all"             # all | bottom | top | buried | [z0, z1]
f = 0.0
\end{verbatim}
Because positions are physical, the same file rasterises on any mesh; the
catalogue shapes written as files reproduce the catalogue masks bit for
bit. The browser editor (\code{julia julia/editor/server.jl}) draws these
files. One picks a face of the sample and paints from it inward, a given
number of cells deep and starting a given number of cells inside, so cuts
can be buried; figures, gradients of $f$ or of depth, and a palette of
condensation strengths complete it. It shows the current, the field and the
face through which vortices enter for each sign of $J$, and checks, while
drawing, the mirror symmetry that would force $\eta=0$, contact with the
injection faces and normal shunts.

% ======================================================================
\section{Protocol for reporting an efficiency}

A critical current is a bracket: the last silent rung and the first
dissipating one of an ascending ladder. A rung integrated for $\Delta t$ at
voltage $V$ carries $|V|\Delta t/2\pi$ vortex crossings per measuring line:
below one it is silent, at ten or more dissipating, and in between undecided
until it is re-measured on a doubled window. A threshold on the voltage
alone, $|V(J)-V(0)|>10^{-4}$, is weaker: at $\Delta t=2000\,t_{\rm GL}$ it
reads a thirtieth of a crossing. The run summary and the terminal menu (key
\code{v}) give both brackets with this rule, $\eta$ at the bracket centres
with its range, the undecided rungs, and the voltage non-reciprocity
$\sigma_V=(|V_+|-|V_-|)/(|V_+|+|V_-|)$ at the first $|J|$ where both
polarities dissipate; a simpler interpolated threshold is printed next to
them for comparison. Before reporting $\eta$:
\begin{enumerate}
\item no region with $0<f<H$, and $V(J=0)$ at the floor;
\item the odd part $[V(+J)-V(-J)]/2$ is nonzero (the drive enters); the
  even part is the rectification;
\item vortices are present somewhere on the ladder;
\item the \code{\_REF} geometry inverts the sign and keeps $|I_c|$;
\item a rigid shift of the defects by $0.05L$ along $y$ leaves the brackets;
\item the asymmetry exceeds that of an $M_y$-symmetric control run with the
  same numerics while it dissipates.
\end{enumerate}
The failure modes these checks catch are described in the methods paper:
an inert injection, a defect touching a contact, normal shunts, index
ranges or mirrors that are not exact, a static drop read as a voltage,
averaging before unwrapping, windows too short near the threshold,
thresholds set by the current grid, floors counted in cycles, fixed-point
tests applied to limit cycles, path dependence of a ladder, a coherence
length other than the intended one, an extinct $\pi$ band still counted,
coefficients acting on the wrong quantity, and contact heat counted as
sample heat.

% ======================================================================
\section{Numerical modes and reproducibility}

\code{Strict()} routes $|z|$, $e^{-i\phi}$ and the complex product through
the C library functions (\code{hypot}, \code{sincos}, a fused multiply--add
product); it is the mode of the regression tests, which compare against
stored states. \code{Fast()} uses Julia's native arithmetic, vectorises,
is about 1.3 times faster and agrees to machine precision. With
\code{use\_threads=true} the planes in $z$ are updated in parallel; each
plane writes distinct cells, so the result does not depend on the number of
threads. At the start of each run a self-test
compares the vectorised kernel with a reference triple loop.

On one core of a laptop (Intel i5-1235U) a cycle of the
$40\times24\times6$ production box costs $0.030$\,s, so a rung on a
$1500\,\tGL$ floor takes about four minutes; with the scalar potential,
$0.046$\,s.

% ======================================================================
\section{Running}

\subsection{From Julia}
\begin{verbatim}
using Pkg; Pkg.activate("julia")
using TDGL
sde = run_case(Config(GEOM_MODE="WEDGE", Nx_sc=40, Ny_sc=24, Nz_sc=6,
                      dx=0.5, dy=0.5, dz=0.5, kappa_s=2.0, xi_ratio=1.30,
                      H_const=0.50, J_list=[0.040, 0.045],
                      t_min_GL=1500.0, t_min_GL_res=2000.0, eps_conv=1e-5,
                      out_dir="runs/first", run_tag="first"))
brackets(sde)        # Ic+- as brackets with the verdict rule, eta and its range
\end{verbatim}
The same, smaller and in about a minute, is \code{examples/first\_run.jl}.

\subsection{From the command line and the terminal menu}

\code{julia/scripts/run\_preset.jl} launches one run with flags for the
geometry (\code{--geom}, \code{--shape}), the parameters (\code{--kappa},
\code{--H}, \code{--xi-ratio}, \code{--gamma12}, \code{--T}), the currents
(\code{--J}), the mesh (\code{--mesh} or \code{--Nx --Ny --Nz --dx}), the
times (\code{--window} or \code{--tmin --tminres}), the scalar potential
(\code{--phi off|coulomb|tf}) and one optional term (\code{--drag},
\code{--heat}, \code{--qss}), and the scheme (\code{--scheme prod|legacy};
prod unless $\Phi$ or $Q$ needs legacy; with heat, drag or $T$, prod only where
its step is at least twice the explicit one; a resumed run keeps its
scheme). \code{julia julia/scripts/tdgl\_tui.jl} builds
the same command step by step, lists the runs, stops them cleanly and
resumes them.

\subsection{Where runs are written}

\code{run\_case} writes to \code{out\_dir}, any path, created if needed;
\code{run\_preset.jl} writes to \code{julia/runs/<name>/} unless \code{--dir}
says otherwise. The terminal menu lists every run below \code{julia/runs}, at
any depth (\code{julia/runs/project/runs/case} is found), and every
\code{runs/} folder one or two levels below the root of the repository
(\code{project/runs}, \code{group/project/runs}); a run kept anywhere else is
not scanned. Runs started outside the menu (a script, a cluster job) appear
as read-only rows, with their progress read from their own output.

\subsection{On a cluster}

The menu is for a workstation; on a cluster a run is one batch job with one
core (\code{JULIA\_NUM\_THREADS=1}, \code{OPENBLAS\_NUM\_THREADS=1}: more BLAS
threads only oversubscribe). Ladders at different fields, shapes or
parameters do not talk to each other, so the natural use of a cluster is a
job array, one task per run: $N$ tasks give $N$ times the throughput with no
change to the solver.
\code{examples/slurm\_array.sh} runs the notched box at five fields:
\begin{verbatim}
#SBATCH --array=0-4 --cpus-per-task=1 --mem=2G --time=24:00:00
export JULIA_NUM_THREADS=1 OPENBLAS_NUM_THREADS=1
FIELDS=(0.015 0.03 0.06 0.09 0.12)
H=${FIELDS[$SLURM_ARRAY_TASK_ID]}
julia --project=julia julia/scripts/run_preset.jl --geom WEDGE --kappa 25 \
      --xi-ratio 4 --H $H --J 0.02,0.04,0.06,0.08,0.10,0.20,0.30,0.40 \
      --mesh production --window long --dir runs/field_H$H ${RESUME:+--restart}
\end{verbatim}
A task stopped by the time limit resumes from its last closed point when it is
submitted again with \code{RESUME=1}: the folder name must not depend on the
job number. \code{examples/slurm\_job.sh} is the same for a single run.
Results are read without the menu: \code{julia julia/scripts/tdgl\_tui.jl
--results}, or \code{brackets(sde)} in Julia.

\subsection{What a run writes}
\begin{center}
\begin{tabular}{ll}
\toprule
file & content \\
\midrule
\code{job.json} & the parameters of the launch \\
\code{j\_history.csv} & one line per closed rung: $J$, $V$, class, vortices \\
\code{run\_info.txt} & time step and cycle length, to count crossings \\
\code{output/RUN\_SUMMARY\_*.txt} & parameters, $I_c^\pm$, $\eta$ and a convergence table \\
\code{restart/state\_J\{pos,neg\}.jld2} & full state (compressed JLD2; \code{.mat} with \code{checkpoint\_format="octave"}) \\
\code{output/snapshots/} & mid-plane $|\Psi|^2$ as PGM and CSV \\
\code{output/vtk/} & legacy VTK for ParaView (\code{vtk\_every}) \\
\bottomrule
\end{tabular}
\end{center}
Any checkpoint converts to VTK with \code{julia/scripts/vtk\_from\_checkpoint.jl}.

% ======================================================================
\section{Tests}

\code{julia --project=julia julia/test/run.jl} runs about $56\,000$
assertions in a minute and a half. They check the kernel against the
reference loop and across threads; regression against two stored states
(a Meissner state and a state with vortex flow);
the lattice Landau level ($1.6\,\%$ low at $\mathrm dx=0.5$, $H=0.5$) and
the surface-nucleation ratio $\Hc3/\Hc2$ ($0.05\,\%$ from $1.69461$);
that $\Gamma_r$ moves no static quantity; that $\Phi$ is a gauge freedom,
also with a Thomas--Fermi length, as long as the residual does not relax;
that the separable solver equals the LU; every optional term and its guards;
shapes read from files; and the VTK writer.

% ======================================================================
\section{Limitations}

\begin{itemize}
\item Standard TDGL with a scalar relaxation time per band; no generalised
  TDGL, no mass tensor, no microscopic derivation of $\Gamma_r$.
\item The two-component functional has two parameters, $m_r$ and $a_{2n}$,
  for three targets of a material such as MgB$_2$ (coherence-length ratio,
  stiffness partition and gap ratio); not all can be matched at once.
\item A notch ($f=0$) is locally at $T_c$: neither vacuum nor a normal metal.
\item The buffer is a normal metal and a systematic: absolute thresholds depend on its width. The
  box has no common electrodes at its ends, so a thick buffer carries normal current in parallel with
  a silent sample (with $24$ cells around the $40\times24\times6$ box, $62\,\%$ of it); one cell keeps
  that path small. A thicker buffer is not a closer model of a film in vacuum.
\item The explicit scheme costs time steps quadratically in $\kappa$; the
  production scheme removes it except with the scalar potential.
\item The silent state is driven, not an equilibrium. The current enters
  through the contact faces as normal current and turns into supercurrent
  within one or two $\xi_\sigma$; those layers carry a steady electric field.
  The current at which the silent state starts to flow therefore depends on
  the relaxation rates and on the contact, not only on the free energy.
\end{itemize}

\end{document}
